\par\Needspace{6\baselineskip}\medskip
\begin{center}\hypertarget{native-185-article138}{}\pdfbookmark[2]{138.}{article-138}138.\end{center}
The \emph{third} task is now to join the first and last geocentric positions of the celestial
body, i.\,e. the points $B,B''$, by a great circle and determine its intersection with the great
circle $A'B'$. Let $B^*$ be this intersection, $\delta'-\sigma$ its distance from the point $A'$, $\alpha^*$ its
longitude, and $\beta^*$ its latitude. Since $B,B^*,B''$ lie on the same great circle, we have the
familiar equation
\[
\begin{aligned}
0={}&\operatorname{tang}\beta\sin(\alpha''-\alpha^*)-\operatorname{tang}\beta^*\sin(\alpha''-\alpha)\\
&+\operatorname{tang}\beta''\sin(\alpha^*-\alpha),
\end{aligned}
\]
which, on substituting $\operatorname{tang}\gamma'\sin(\alpha^*-l')$ for $\operatorname{tang}\beta^*$, takes the following form
\[
\begin{aligned}
0={}&\cos(\alpha^*-l')\{\operatorname{tang}\beta\sin(\alpha''-l')-\operatorname{tang}\beta''\sin(\alpha-l')\}\\
&-\sin(\alpha^*-l')\{\operatorname{tang}\beta\cos(\alpha''-l')\\
&\qquad+\operatorname{tang}\gamma'\sin(\alpha''-\alpha)-\operatorname{tang}\beta''\cos(\alpha-l')\}.
\end{aligned}
\]
Hence, since $\operatorname{tang}(\alpha^*-l')=\cos\gamma'\operatorname{tang}(\delta'-\sigma)$, we have
\[
\begin{gathered}
\operatorname{tang}(\delta'-\sigma)=\\[.4ex]
\frac{\operatorname{tang}\beta\sin(\alpha''-l')-\operatorname{tang}\beta''\sin(\alpha-l')}{\begin{gathered}\cos\gamma'(\operatorname{tang}\beta\cos(\alpha''-l')-\operatorname{tang}\beta''\cos(\alpha-l'))\\+\sin\gamma'\sin(\alpha''-\alpha)\end{gathered}}\mathpunct{\cdot}
\end{gathered}
\]

From this we derive the following formulas, which are better suited to numerical calculation.
Set
\[
\begin{aligned}
7)\quad&\operatorname{tang}\beta\sin(\alpha''-l')-\operatorname{tang}\beta''\sin(\alpha-l')=S\\[.5ex]
8)\quad&\operatorname{tang}\beta\cos(\alpha''-l')-\operatorname{tang}\beta''\cos(\alpha-l')=T\sin t\\[.5ex]
9)\quad&\sin(\alpha''-\alpha)=T\cos t
\end{aligned}
\]
(Article 14, II) to obtain
\[
10)\quad\operatorname{tang}(\delta'-\sigma)=\frac{S}{T\sin(t+\gamma')}\mathpunct{\cdot}
\]
The ambiguity in determining the arc $\delta'-\sigma$ from its tangent arises because the great circles
$A'B',BB''$ intersect at \emph{two} points. For $B^*$ we shall always \hypertarget{native-186}{}choose the intersection
nearest to the point $B'$, so that $\sigma$ always lies between the limits $-90^0$ and $+90^0$.
This removes the ambiguity. The value of the arc $\sigma$ (which depends on the \emph{curvature}
of the geocentric motion) will then usually be fairly small, and in general will be of the
second order if the time intervals are regarded as quantities of the first order.

The changes to the calculation when the equator is chosen as the fundamental plane in place of
the ecliptic follow directly from the remark in the preceding article.

It is also clear that the position of the point $B^*$ would remain indeterminate if the circles
$BB'',A'B'$ coincided altogether. We exclude this case, in which the four points $A',B,B',B''$ would lie on
the same great circle, from our investigation. When choosing observations, it is also advisable
to avoid cases in which the four points lie close to a single great circle. In such a case even
the smallest errors in the observations would have too great an effect on the position of the
point $B^*$, which is important in the operations that follow, and its position could not be
determined with the required precision. --- It is likewise clear that the point $B^*$ remains
indeterminate whenever the points $B,B''$ coincide\footnote{Or when they are opposite each other.
We do not discuss this case because our method should not be extended to observations spanning
so long an interval.}, in which case the orientation of the circle $BB''$ itself would be
indeterminate. We therefore also exclude this case. For reasons similar to those just given, we
must also avoid observations in which the first and last geocentric positions fall at points of
the sphere close to each other.
